\documentclass[11pt,a4paper]{article}

% ---------- Pacchetti ----------
\usepackage[italian]{babel}
\usepackage[utf8]{inputenc}
\usepackage[T1]{fontenc}
\usepackage[margin=2.5cm]{geometry}
\usepackage{graphicx}
\usepackage{tikz}
\usetikzlibrary{arrows.meta,positioning,fit,backgrounds,shapes.multipart,calc}
\usepackage{booktabs}
\usepackage{multirow}
\usepackage{xcolor}
\usepackage{hyperref}
\usepackage{listings}
\usepackage{caption}
\usepackage{enumitem}
\usepackage{parskip}

\hypersetup{
    colorlinks=true,
    linkcolor=blue!50!black,
    urlcolor=blue!50!black
}

% ---------- Colori ----------
\definecolor{doColor}{RGB}{224,242,254}
\definecolor{doBorder}{RGB}{2,132,199}
\definecolor{luColor}{RGB}{254,242,224}
\definecolor{luBorder}{RGB}{194,120,3}
\definecolor{brokerColor}{RGB}{237,233,254}
\definecolor{brokerBorder}{RGB}{109,40,217}
\definecolor{dbColor}{RGB}{220,252,231}
\definecolor{dbBorder}{RGB}{21,128,61}
\definecolor{clusterA}{RGB}{219,234,254}
\definecolor{clusterB}{RGB}{254,226,226}
\definecolor{clusterC}{RGB}{220,252,231}

\lstdefinestyle{jsonstyle}{
    basicstyle=\ttfamily\small,
    breaklines=true,
    columns=fullflexible,
    backgroundcolor=\color{gray!8},
    frame=single,
    rulecolor=\color{gray!40},
    xleftmargin=4pt,
    xrightmargin=4pt
}
\lstset{style=jsonstyle}

\title{\textbf{Come la Linkage Unit costruisce, persiste \\ e aggiorna la Link Table}\\[4pt]
\large Panoramica architetturale ad alto livello --- progetto PRIMAT}
\author{}
\date{}

\begin{document}
\maketitle
\vspace{-1.2cm}

\section{Introduzione}

PRIMAT \`e un sistema di \emph{Privacy-Preserving Record Linkage} (PPRL): pi\`u
\textbf{Data Owner} (DO) --- ciascuno titolare di una propria base dati contenente
informazioni identificative su persone --- vogliono scoprire quali dei loro record
si riferiscono alla stessa entit\`a reale, senza mai condividere i dati in chiaro
tra loro. A occuparsi di questo compito \`e la \textbf{Linkage Unit} (LU): un
componente centrale che riceve dai Data Owner solo impronte offuscate dei record
(\emph{Bloom Filter}), calcola quali impronte si assomigliano abbastanza da
rappresentare la stessa persona, e mantiene nel tempo una \textbf{Link Table}:
la struttura che collega i record delle varie sorgenti allo stesso identificativo
pseudonimizzato.

Questo documento descrive, ad alto livello e senza entrare nei dettagli
implementativi o matematici, quattro aspetti:
\begin{enumerate}[nosep]
    \item come la LU costruisce la Link Table la prima volta;
    \item come questa struttura viene resa persistente;
    \item come funziona l'aggiornamento incrementale quando arrivano nuovi record;
    \item come l'ambiente cos\`i prodotto si presta naturalmente a essere interrogato
        da un motore di query federata (es. Trino) che unisce le sorgenti in una
        vista globale pseudonimizzata.
\end{enumerate}

\section{Costruzione iniziale della Link Table}

\subsection{Il flusso concettuale}

Quando la Linkage Unit viene avviata per la prima volta su un insieme di Data
Owner, non esiste ancora alcuna Link Table: viene costruita da zero seguendo
quattro fasi concettuali, applicate a tutte le impronte ricevute:

\begin{enumerate}
    \item \textbf{Blocking}: invece di confrontare ogni record con tutti gli
        altri (costo quadratico), i record vengono raggruppati in "blocchi"
        tramite chiavi calcolate con tecniche di \emph{Locality Sensitive
        Hashing} (LSH). Solo i record che condividono almeno una chiave di
        blocco vengono confrontati fra loro.
    \item \textbf{Calcolo della similarit\`a}: per ogni coppia di record
        candidata si calcola un punteggio di somiglianza tra le rispettive
        impronte.
    \item \textbf{Classificazione}: le coppie con similarit\`a sopra una soglia
        configurabile vengono marcate come "match" (stessa entit\`a probabile).
    \item \textbf{Clustering}: i match vengono aggregati in gruppi (cluster) di
        record che rappresentano la stessa persona, anche quando la catena di
        somiglianze attraversa pi\`u di due sorgenti.
    \end{enumerate}

Il risultato di questa pipeline \`e proprio la Link Table iniziale: un insieme
di cluster, ciascuno identificato da un id univoco (lo pseudonimo globale), che
raggruppa i record provenienti da una o pi\`u sorgenti.

\subsection{Configurabilit\`a}

L'intero processo \`e guidato da un file di configurazione dichiarativo, non da
codice: si specificano le sorgenti (Data Owner) coinvolte, la soglia di
similarit\`a, i parametri dell'algoritmo di blocking e, soprattutto, la
\textbf{strategia di clustering} da applicare. Sono disponibili pi\`u strategie
intercambiabili, ciascuna con i propri parametri di tuning, pensate per scenari
diversi (maggiore tolleranza al rumore, maggiore velocit\`a, garanzie diverse
sulla qualit\`a dei cluster). La scelta della strategia e dei suoi parametri
avviene quindi semplicemente cambiando la configurazione, senza toccare la
logica applicativa.

\subsection{Supporto di tutti i contesti dirty/clean}

Una sorgente pu\`o essere \textbf{"dirty"} (pu\`o contenere pi\`u record che si
riferiscono alla stessa persona, cio\`e ha bisogno anche di una deduplica
interna) oppure \textbf{"clean"} (gi\`a deduplicata al suo interno). La LU
supporta in modo nativo tutte le combinazioni:

\begin{itemize}[nosep]
    \item \textbf{dirty--dirty}: nessuna sorgente \`e gi\`a deduplicata;
    \item \textbf{dirty--clean (mista)}: una o pi\`u sorgenti pulite convivono
        con sorgenti sporche;
    \item \textbf{clean--clean}: tutte le sorgenti sono gi\`a deduplicate,
        serve solo collegarle tra loro;
    \item \textbf{sorgente singola (deduplica)}: un solo Data Owner "dirty"
        vuole scoprire i duplicati interni alla propria base dati, senza alcun
        secondo Data Owner coinvolto.
\end{itemize}

Ognuno di questi scenari \`e semplicemente una diversa combinazione di
configurazione (quali sorgenti sono dichiarate "pulite" e quale strategia di
clustering si sceglie), non un percorso di codice separato.

\subsection{Multi-sorgente per natura}

Il matching non \`e limitato a due sole sorgenti confrontate a coppie: la LU
costruisce un unico grafo di similarit\`a che pu\`o coinvolgere
contemporaneamente un numero arbitrario di Data Owner. Un cluster pu\`o quindi
raggruppare record provenienti da tre, quattro o pi\`u sorgenti differenti in
un solo passaggio, se la catena di somiglianze lo giustifica.

\subsection{Relazione tra Data Owner e Linkage Unit}

I Data Owner non inviano mai dati identificativi in chiaro. Ogni Data Owner
trasforma localmente i propri dati in un'impronta offuscata e irreversibile
(un Bloom Filter costruito a partire da frammenti hashati dei campi
identificativi) e la pubblica su un canale di messaggistica intermedio; la
Linkage Unit si mette in ascolto su quello stesso canale, raccoglie le
impronte di tutte le sorgenti coinvolte e avvia il calcolo. Il canale funge
anche da tramite per i comandi di avvio inviati dalla LU ai Data Owner.

\begin{figure}[h!]
\centering
\begin{tikzpicture}[
    node distance=1.4cm and 2.6cm,
    box/.style={rectangle, rounded corners, draw, minimum width=3cm, minimum height=1.1cm, align=center, font=\small},
    do/.style={box, fill=doColor, draw=doBorder, thick},
    broker/.style={box, fill=brokerColor, draw=brokerBorder, thick, minimum width=3.4cm},
    lu/.style={box, fill=luColor, draw=luBorder, thick},
    lbl/.style={font=\scriptsize, align=center},
    arr/.style={-{Stealth[length=2.2mm]}, thick}
]
    \node[do] (doa) {Data Owner A\\[1pt]\scriptsize (dirty)};
    \node[do, below=of doa] (dob) {Data Owner B\\[1pt]\scriptsize (clean)};
    \node[do, below=of dob] (doc) {Data Owner C\\[1pt]\scriptsize (dirty)};

    \node[broker, right=of dob] (broker) {Broker MQTT\\[1pt]\scriptsize (Moquette)};

    \node[lu, right=of broker] (lu) {Linkage Unit};

    \draw[arr] (doa) to[bend left=12] node[lbl, above, sloped]{\texttt{rbf/A}} (broker);
    \draw[arr] (dob) -- node[lbl, above]{\texttt{rbf/B}} (broker);
    \draw[arr] (doc) to[bend right=12] node[lbl, below, sloped]{\texttt{rbf/C}} (broker);

    \draw[arr] (broker) -- node[lbl, above]{RBF ricevuti} (lu);

    \draw[arr] (lu) to[bend left=18] node[lbl, below]{\texttt{cmd} (avvio run)} (broker);
    \draw[arr] (broker) to[bend left=12] (doa);
    \draw[arr] (broker) to[bend right=12] (doc);

    \node[lbl, below=0.35cm of doc, xshift=2.2cm, text width=8.2cm] {
        Sul canale viaggiano \textbf{solo impronte offuscate} (Bloom Filter) e
        digest di configurazione: mai valori identificativi in chiaro.
    };
\end{tikzpicture}
\caption{Relazione tra Data Owner e Linkage Unit: ciascun Data Owner pubblica
la propria impronta offuscata su un broker di messaggistica condiviso; la
Linkage Unit riceve le impronte di tutte le sorgenti coinvolte e coordina
l'avvio di ogni esecuzione.}
\end{figure}

\clearpage
\section{Persistenza della Link Table}

\subsection{Il modello di storage}

La Link Table non \`e un file volatile ricalcolato da zero a ogni esecuzione:
viene salvata in un database relazionale, cos\`i che le esecuzioni successive
possano ripartire da quanto gi\`a consolidato. Concettualmente lo schema \`e
composto da poche entit\`a:

\begin{itemize}[nosep]
    \item \textbf{Party}: rappresenta un Data Owner registrato presso la LU.
    \item \textbf{Record}: l'impronta di un singolo record ricevuto da una
        Party, con un identificativo locale univoco (prefissato per evitare
        collisioni tra sorgenti diverse).
    \item \textbf{Cluster}: l'entit\`a "reale" riconosciuta nel tempo --- \`e
        di fatto la riga della Link Table: il suo id \`e lo pseudonimo globale
        condiviso da tutti i record che raggruppa.
    \item \textbf{Chiavi di blocco}: le stesse chiavi usate in fase di
        blocking vengono conservate accanto a record e cluster, cos\`i da
        poter individuare rapidamente, in futuro, quali cluster storici sono
        potenzialmente coinvolti da un nuovo record.
\end{itemize}

Ogni strategia di clustering scrive su un proprio spazio dati dedicato: questo
mantiene isolate esecuzioni con strategie diverse e permette di confrontarne
i risultati senza interferenze reciproche.

\begin{figure}[h!]
\centering
\begin{tikzpicture}[
    node distance=1.6cm and 2.2cm,
    ent/.style={rectangle, rounded corners, draw, thick, align=center, font=\small, minimum height=1cm},
    party/.style={ent, fill=doColor, draw=doBorder, minimum width=2.6cm},
    rec/.style={ent, fill=luColor, draw=luBorder, minimum width=2.6cm},
    clu/.style={ent, fill=dbColor, draw=dbBorder, minimum width=2.6cm},
    blk/.style={ent, fill=brokerColor, draw=brokerBorder, minimum width=3.1cm},
    rel/.style={font=\scriptsize},
    line/.style={thick}
]
    \node[party] (party) {Party};
    \node[rec, right=of party] (record) {Record};
    \node[clu, right=of record] (cluster) {Cluster};
    \node[blk, below=1.7cm of record, xshift=1.1cm] (block) {Chiave di blocco\\[1pt]\scriptsize (BlockingKeyAttribute)};

    \draw[line] (party) -- (record) node[rel, midway, above]{1 -- N};
    \draw[line] (record) -- (cluster) node[rel, midway, above]{N -- 1};
    \draw[line] (record) to[bend right=15] node[rel, left]{N -- M} (block);
    \draw[line] (cluster) to[bend left=15] node[rel, right]{N -- M} (block);

    \node[ent, dashed, below=1.6cm of party, minimum width=3.6cm] (state)
        {PartyEncodingState\\[1pt]\scriptsize (config attesa per party)};
    \draw[line, dashed] (party) -- (state) node[rel, midway, left]{1 -- 1};
\end{tikzpicture}
\caption{Schema concettuale della persistenza: ogni Party ha molti Record, ogni
Record appartiene a un solo Cluster (finch\'e non viene eventualmente fuso),
Record e Cluster condividono le chiavi di blocco che permettono di ritrovarli
rapidamente in futuro. \texttt{PartyEncodingState} verifica che una Party non
cambi silenziosamente il modo in cui genera le proprie impronte tra
un'esecuzione e l'altra.}
\end{figure}

\clearpage
\subsection{Esempio minimo: prima inizializzazione}

Supponiamo due Data Owner, \textbf{A} (dirty) e \textbf{B} (clean), che inviano
alla LU le impronte di pochi record. Per rendere l'esempio leggibile, mostriamo
qui i valori identificativi originali dei record: nella realt\`a questi valori
\emph{non lasciano mai} il Data Owner, che trasmette solo l'impronta offuscata
corrispondente.

\begin{table}[h!]
\centering
\caption{Record in ingresso al primo run (valori mostrati solo a scopo
illustrativo)}
\begin{tabular}{@{}llll@{}}
\toprule
\textbf{Id record} & \textbf{Party} & \textbf{Nome} & \textbf{Cognome} \\
\midrule
A1 & A & Mario   & Rossi \\
A2 & A & Mario   & Rossi \\
A3 & A & Luca    & Bianchi \\
B1 & B & Maria   & Verdi \\
B2 & B & Mario   & Rossi \\
\bottomrule
\end{tabular}
\end{table}

Dopo blocking, calcolo di similarit\`a, classificazione e clustering, la LU
riconosce che \texttt{A1}, \texttt{A2} e \texttt{B2} rappresentano la stessa
persona (Mario Rossi, con un duplicato interno alla sorgente A), mentre
\texttt{A3} e \texttt{B1} restano singoli. La Link Table risultante:

\begin{table}[h!]
\centering
\caption{Link Table dopo la prima inizializzazione}
\begin{tabular}{@{}lll@{}}
\toprule
\textbf{Cluster (pseudonimo globale)} & \textbf{Record contenuti} & \textbf{Party coinvolte} \\
\midrule
Cluster 1 & A1, A2, B2 & A, B \\
Cluster 2 & A3         & A \\
Cluster 3 & B1         & B \\
\bottomrule
\end{tabular}
\end{table}

\begin{figure}[h!]
\centering
\begin{tikzpicture}[
    chip/.style={rectangle, rounded corners, draw, thick, minimum width=1.3cm, minimum height=0.8cm, font=\scriptsize, align=center},
    lbl/.style={font=\scriptsize\itshape}
]
    \node[chip, fill=doColor, draw=doBorder] (a1) at (0,0) {A1};
    \node[chip, fill=doColor, draw=doBorder, right=0.35cm of a1] (a2) {A2};
    \node[chip, fill=luColor, draw=luBorder, right=0.35cm of a2] (b2) {B2};

    \node[chip, fill=doColor, draw=doBorder, right=2.1cm of b2] (a3) {A3};

    \node[chip, fill=luColor, draw=luBorder, right=2.1cm of a3] (b1) {B1};

    \begin{scope}[on background layer]
        \node[fit=(a1)(a2)(b2), draw=doBorder, thick, rounded corners,
              fill=clusterA!35, inner sep=9pt,
              label={above:{\bfseries\small Cluster 1}}] {};
        \node[fit=(a3), draw=doBorder, thick, rounded corners,
              fill=clusterB!45, inner sep=9pt,
              label={above:{\bfseries\small Cluster 2}}] {};
        \node[fit=(b1), draw=luBorder, thick, rounded corners,
              fill=clusterC!45, inner sep=9pt,
              label={above:{\bfseries\small Cluster 3}}] {};
    \end{scope}

    \node[lbl, below=0.9cm of a3, xshift=0cm, text width=10.5cm] {
        L'id di ciascun cluster \`e lo pseudonimo globale usato per collegare
        record di sorgenti diverse senza mai esporre i dati originali.
    };
\end{tikzpicture}
\caption{Vista concettuale della Link Table: ogni cluster \`e un contenitore
pseudonimizzato che raggruppa i record (di una o pi\`u sorgenti) riconosciuti
come la stessa entit\`a.}
\end{figure}

\clearpage
\section{Logica di Incrementing}

Una volta che la Link Table esiste, la LU non deve ripartire da zero ogni
volta che arrivano nuovi record: li integra in modo incrementale.

\subsection{Il principio}

Quando arriva un nuovo lotto di record:

\begin{enumerate}
    \item la LU non ricarica l'intero storico, ma recupera \emph{solo} i
        cluster gi\`a esistenti che condividono almeno una chiave di blocco
        con uno dei record nuovi (i "candidati");
    \item l'appartenenza ai cluster gi\`a persistiti viene presa come punto di
        partenza fisso: un legame stabilito in un'esecuzione precedente non
        pu\`o mai essere smontato da un'esecuzione successiva;
    \item i nuovi match vengono poi applicati sopra questa base, e per ogni
        gruppo di record risultante si determina l'esito confrontandolo con i
        cluster storici coinvolti:
    \begin{itemize}[nosep]
        \item \textbf{nessun cluster storico coinvolto} $\rightarrow$ si crea
            un \emph{nuovo cluster};
        \item \textbf{un solo cluster storico coinvolto} $\rightarrow$ il
            cluster \emph{cresce}, incorporando i record nuovi;
        \item \textbf{pi\`u cluster storici coinvolti} $\rightarrow$ i cluster
            vengono \emph{fusi} in uno solo (l'id pi\`u basso viene mantenuto
            come identificativo del cluster risultante).
    \end{itemize}
\end{enumerate}

In altre parole: un cluster pu\`o solo crescere o fondersi con altri, mai
essere diviso silenziosamente da un'esecuzione successiva.

\subsection{Esempio minimo: arrivo di nuovi record}

Riprendiamo la Link Table dell'esempio precedente (Cluster 1 = \{A1, A2, B2\},
Cluster 2 = \{A3\}, Cluster 3 = \{B1\}).

\medskip
\noindent\textbf{Caso 1 --- crescita di un cluster.} Un terzo Data Owner
\textbf{C} invia un nuovo record \texttt{C1 = "Mario Rossi"}. La LU trova che
\texttt{C1} corrisponde al Cluster 1 gi\`a esistente:

\begin{table}[h!]
\centering
\begin{tabular}{@{}lll@{}}
\toprule
 & \textbf{Prima} & \textbf{Dopo l'arrivo di C1} \\
\midrule
Cluster 1 & A1, A2, B2 & A1, A2, B2, \textbf{C1} \\
\bottomrule
\end{tabular}
\end{table}

\noindent\textbf{Caso 2 --- fusione di due cluster.} Successivamente arriva un
nuovo record \texttt{C2} da C che, a sua volta, risulta simile sia a
\texttt{A3} (Cluster 2) sia a \texttt{B1} (Cluster 3): \`e un "record ponte"
che rivela che i due cluster storici in realt\`a rappresentano la stessa
persona.

\begin{table}[h!]
\centering
\begin{tabular}{@{}lll@{}}
\toprule
 & \textbf{Prima} & \textbf{Dopo l'arrivo di C2} \\
\midrule
Cluster 2 & A3  & \multirow{2}{*}{Cluster 2 (fuso): A3, B1, \textbf{C2}} \\
Cluster 3 & B1  &  \\
\bottomrule
\end{tabular}
\end{table}

Il Cluster 3 viene svuotato e i suoi record vengono ripuntati sul Cluster 2
(id pi\`u basso), che diventa il cluster risultante della fusione.

\section{Estensioni possibili: DELETE e UPDATE di record storici}

Lo schema descritto finora copre l'inserimento di record nuovi, ma non ancora
la loro rimozione o correzione. \`E ragionevole immaginare un'estensione in
cui un Data Owner possa richiedere esplicitamente, tramite un canale
dedicato, la cancellazione o la modifica di un proprio record gi\`a linkato
in passato. Si tratta di una funzionalit\`a \textbf{non implementata a oggi},
qui discussa solo a livello di principio:

\begin{itemize}
    \item \textbf{DELETE esplicito}: il Data Owner segnala l'id di un proprio
        record da rimuovere. La LU dovrebbe scollegare il record dal proprio
        cluster senza dover ripetere l'intero calcolo di linkage: se il
        cluster conteneva altri record, resta valido e si riduce soltanto; se
        il record era l'unico membro del cluster, il cluster stesso viene
        eliminato.
    \item \textbf{UPDATE esplicito}: una correzione ai dati di un record pu\`o
        essere trattata come una cancellazione seguita da un nuovo inserimento
        con l'impronta aggiornata, rientrando cos\`i nella normale logica di
        incrementing gi\`a descritta, senza bisogno di un percorso separato.
    \item \textbf{Implicazioni da considerare}: un canale di richiesta
        esplicito e autenticato lato Data Owner, tracciabilit\`a delle
        modifiche per finalit\`a di audit, e la possibilit\`a che la rimozione
        di un record "ponte" renda necessario valutare se un cluster fuso in
        passato debba essere ricontrollato.
\end{itemize}

\section{Verso l'interrogazione federata: come la Link Table prepara il terreno}

L'integrazione di un motore di query federata (ad esempio Trino) che unisca le
sorgenti dei Data Owner e la Link Table in un'unica vista globale \`e in corso
di sviluppo separatamente e non \`e oggetto di questo documento. Vale per\`o la
pena osservare \emph{perch\'e} l'ambiente che la Linkage Unit gi\`a produce si
presta naturalmente a questo tipo di interrogazione:

\begin{itemize}
    \item \textbf{Schema relazionale standard}: la Link Table vive in un
        database relazionale interrogabile con un connettore JDBC comune, senza
        bisogno di adattatori dedicati o formati proprietari.
    \item \textbf{Chiave di join gi\`a pseudonimizzata}: l'id di cluster \`e
        esattamente lo pseudonimo globale che un motore federato userebbe come
        chiave di join tra le tabelle dei diversi Data Owner, senza mai dover
        transitare per dati identificativi in chiaro.
    \item \textbf{Identificativi senza collisioni}: gli id dei record sono
        prefissati per sorgente, cos\`i che una vista che unisce pi\`u Data
        Owner non incontri mai ambiguit\`a tra record di provenienza diversa.
    \item \textbf{Spazi dati isolati e prevedibili}: ogni strategia di
        clustering scrive nel proprio spazio dati dedicato, cos\`i un motore
        di interrogazione pu\`o scegliere in modo esplicito quale "versione"
        della Link Table consultare, senza risultati misti o ambigui.
    \item \textbf{Coerenza gi\`a garantita a monte}: i controlli di coerenza
        che la LU esegue tra le sorgenti (stessa configurazione di codifica,
        stessa dimensione delle impronte) assicurano che i dati letti da un
        motore federato siano sempre internamente consistenti, senza bisogno
        che sia il motore di query a doverlo verificare.
\end{itemize}

In sintesi: il lavoro della Linkage Unit non produce soltanto i collegamenti
tra i record, ma un ambiente dati gi\`a pulito, stabile e pseudonimizzato ---
esattamente il presupposto di cui un livello di virtual data integration ha
bisogno per poter interrogare pi\`u sorgenti come se fossero un'unica vista
coerente.

\end{document}
